\section{Đặc tả Usecase Module 1}
% --- //////////////////////////////////////////////////////////////// ---
% MODULE 1

\renewcommand{\arraystretch}{1.2}
\begin{longtable}{@{}p{4cm} p{11cm}@{}}
\caption{Mô tả Use Case: Tạo lộ trình}
\label{usecase-1-01} \\
\toprule
\textbf{Use case ID}       & UC-M1-01 \\ \midrule
\textbf{Tên Use case}      & Tạo lộ trình \\ \midrule
\textbf{Các tác nhân}      & Traveler \\ \midrule
\textbf{Mô tả}             & Khởi tạo một khung sườn lộ trình mới. Traveler có thể chọn tạo một lộ trình rỗng, hoặc tạo một bản sao từ một lộ trình đã có. Traveler cũng sẽ đặt tên cho lộ trình mới này. \\ \midrule
\textbf{Tiền điều kiện}    & Traveler đã đăng nhập và đang ở trang quản lý các hành trình. \\ \midrule
\textbf{Hậu điều kiện}     & Một bản nháp lộ trình mới (rỗng hoặc được sao chép) được tạo ra trong phiên làm việc của người dùng. \\ \midrule
\textbf{Luồng sự kiện chính} & 
\begin{tabular}[t]{@{}p{10.5cm}@{}}
1. Traveler nhấn ""Tạo hành trình mới"".\\
2. Hệ thống hỏi Traveler muốn ""Tạo mới hoàn toàn"" hay ""Tạo từ hành trình đã có"".\\
3. Traveler chọn ""Tạo mới hoàn toàn"".\\
4. Hệ thống yêu cầu đặt tên cho lộ trình.\\
5. Traveler nhập tên (hoặc bỏ trống để lấy tên mặc định).\\
6. Hệ thống tạo ra một lộ trình rỗng với tên đã đặt và hiển thị giao diện chỉnh sửa.
\end{tabular} \\ \midrule
\textbf{Luồng rẽ nhánh} &
\begin{tabular}[t]{@{}p{10.5cm}@{}}
\textbf{A1: Tạo từ hành trình đã có:}\\
3a. Traveler chọn ""Tạo từ hành trình đã có"".\\
4a. Hệ thống hiển thị danh sách các hành trình cũ.\\
5a. Traveler chọn một hành trình.\\
6a. Hệ thống yêu cầu đặt tên cho bản sao mới.\\
7a. Traveler nhập tên.\\
8a. Hệ thống tạo ra một bản sao của hành trình đã chọn.
\end{tabular} \\ \midrule
\textbf{Luồng ngoại lệ} &
\begin{tabular}[t]{@{}p{10.5cm}@{}}
\textbf{E1: Lỗi hệ thống hoặc kết nối mạng:}\\
Luồng này có thể xảy ra tại bước 6 hoặc 8a:\\
6a. Hệ thống không thể tạo đối tượng lộ trình mới do lỗi cơ sở dữ liệu hoặc lỗi kết nối.\\
6b. Hệ thống hiển thị thông báo lỗi chung (ví dụ: ""Không thể tạo hành trình mới lúc này, vui lòng thử lại sau.""). Use case kết thúc.\\[3pt]
\textbf{E2: Tên hành trình không hợp lệ:}\\
Luồng này xảy ra tại bước 5 hoặc 7a:\\
5a. Traveler nhập một tên đã tồn tại (nếu hệ thống yêu cầu tên là duy nhất) hoặc chứa ký tự không hợp lệ.\\
5b. Hệ thống chặn việc tạo và hiển thị thông báo lỗi tương ứng (ví dụ: ""Tên hành trình này đã tồn tại, vui lòng chọn tên khác."").\\[3pt]
\textbf{E3: Người dùng chưa có hành trình cũ:}\\
Luồng này xảy ra tại bước 4a.:\\
4a.1. Hệ thống hiển thị một danh sách rỗng\\
4a.2. Traveler quay lại trang trước đó
\end{tabular} \\ 
\bottomrule
\end{longtable}

\renewcommand{\arraystretch}{1.2}
\begin{longtable}{@{}p{4cm} p{11cm}@{}}
\caption{Mô tả Use Case: Nhập yêu cầu gợi ý lộ trình du lịch}
\label{usecase-1-02} \\
\toprule
\textbf{Use case ID}       & UC-M1-02 \\ \midrule
\textbf{Tên Use case}      & Nhập yêu cầu gợi ý lộ trình du lịch \\ \midrule
\textbf{Các tác nhân}      & Traveler \\ \midrule
\textbf{Mô tả}             & Traveler nhập một yêu cầu (prompt) bằng ngôn ngữ tự nhiên để nhận một lộ trình du lịch do AI tạo ra. \\ \midrule
\textbf{Tiền điều kiện}    & Traveler đang ở trong giao diện chỉnh sửa của một lộ trình rỗng (vừa được tạo từ UC-M1-01). \\ \midrule
\textbf{Hậu điều kiện}     & Một hoặc nhiều bản nháp của lộ trình được hiển thị trên màn hình để Traveler lựa chọn và tinh chỉnh. \\ \midrule
\textbf{Luồng sự kiện chính} &
\begin{tabular}[t]{@{}p{10.5cm}@{}}
1. Traveler nhập một prompt đầy đủ và hợp lệ.\\
2. Hệ thống phân tích và trích xuất thành công các thực thể (yêu cầu, ràng buộc).\\
3. Dựa trên các yêu cầu đó, mô hình AI tạo ra một lộ trình tối ưu và đáp ứng chính xác 100\% các tiêu chí.\\
4. Hệ thống hiển thị lộ trình vừa được AI tạo ra cho Traveler.
\end{tabular} \\ \midrule
\textbf{Luồng rẽ nhánh} &
\begin{tabular}[t]{@{}p{10.5cm}@{}}
\textbf{A1: Không thể tạo lộ trình khớp 100\% yêu cầu:}\\
Tại bước 3 của luồng chính, mô hình AI xác định rằng không thể xây dựng một lộ trình thỏa mãn tất cả các ràng buộc cứng (ví dụ: ngân sách không đủ cho số ngày và loại hình hoạt động yêu cầu):\\
3a . Hệ thống sẽ tự động điều chỉnh (nới lỏng) một vài ràng buộc và tạo ra 1-3 phương án lộ trình thay thế khả thi nhất.\\
3b . Hệ thống hiển thị các lộ trình này, kèm theo giải thích về sự điều chỉnh mà AI đã thực hiện.\\
Use case tiếp tục bình thường.
\end{tabular} \\ \midrule
\textbf{Luồng ngoại lệ} &
\begin{tabular}[t]{@{}p{10.5cm}@{}}
\textbf{E1: Prompt của người dùng không hợp lệ:}\\
Tại bước 1, Traveler nhập một prompt không rõ ràng:\\
1a . Hệ thống không thể trích xuất đủ thông tin để AI có thể bắt đầu quá trình sáng tạo.\\
1b . Hệ thống hiển thị thông báo lỗi, giải thích và yêu cầu Traveler cung cấp 1 prompt chi tiết hơn.
\end{tabular} \\
\bottomrule
\end{longtable}

\renewcommand{\arraystretch}{1.2}
\begin{longtable}{@{}p{4cm} p{11cm}@{}}
\caption{Mô tả Use Case: Đề xuất thông tin Nhà cung cấp Dịch vụ}
\label{usecase-1-03} \\
\toprule
\textbf{Use case ID}       & UC-M1-03 \\ \midrule
\textbf{Tên Use case}      & Đề xuất thông tin Nhà cung cấp Dịch vụ \\ \midrule
\textbf{Các tác nhân}      & Traveler \\ \midrule
\textbf{Mô tả}             & Cho phép Traveler yêu cầu hệ thống cung cấp thông tin liên hệ hoặc thông tin nhà cung cấp cho các địa điểm/dịch vụ đã có sẵn trong lộ trình. Chức năng này giúp người dùng tham khảo và chủ động liên hệ với các nhà cung cấp khi cần thiết. \\ \midrule
\textbf{Tiền điều kiện}    & 
\begin{tabular}[t]{@{}p{10.5cm}@{}}
1. Traveler đã đăng nhập vào hệ thống.\\
2. Traveler đang xem một lộ trình có chứa ít nhất một [Địa điểm] hoặc [Dịch vụ].
\end{tabular} \\ \midrule
\textbf{Hậu điều kiện}     & (Thành công): Hệ thống hiển thị thông tin nhà cung cấp hoặc các dịch vụ liên quan cho mục đã chọn.\\ (Thất bại): Hệ thống thông báo không tìm thấy thông tin. \\ \midrule
\textbf{Luồng sự kiện chính} &
\begin{tabular}[t]{@{}p{10.5cm}@{}}
1. Traveler đang xem chi tiết một lộ trình.\\
2. Traveler chọn một mục trong lộ trình (ví dụ: [Địa điểm] Dinh Độc Lập hoặc [Dịch vụ] Tour 2 đảo).\\
3. Traveler nhấn nút ""Đề xuất thông tin Nhà cung cấp"".\\
4. Hệ thống kiểm tra loại của mục đã chọn.\\
5. (Nếu mục là [Địa điểm]): Hệ thống truy vấn CSDL để tìm các dịch vụ con có đăng ký tại địa điểm đó (ví dụ: ""Dịch vụ thuyết minh"", ""Quầy vé"", ""Xe điện tham quan"") và hiển thị danh sách các dịch vụ này kèm thông tin liên hệ (nếu có).\\
6. (Nếu mục là [Dịch vụ]): Hệ thống truy vấn CSDL để tìm thông tin liên hệ (số điện thoại, website, địa chỉ) của nhà cung cấp dịch vụ đó (ví dụ: ""Thông tin liên hệ của Rooty Trip"").\\
7. Hệ thống hiển thị các thông tin tìm được cho Traveler tham khảo.
\end{tabular} \\ \midrule
\textbf{Luồng rẽ nhánh}    & N/A \\ \midrule
\textbf{Luồng ngoại lệ} &
\begin{tabular}[t]{@{}p{10.5cm}@{}}
\textbf{E1: Không tìm thấy thông tin:}\\
Tại bước 5 hoặc 6, hệ thống không tìm thấy thông tin nhà cung cấp hoặc dịch vụ liên quan nào trong CSDL.\\
Hệ thống hiển thị thông báo ""Không tìm thấy thông tin nhà cung cấp/dịch vụ cho mục này.""\\[3pt]
\textbf{E2: Mục không hỗ trợ:}\\
Tại bước 2, Traveler chọn một mục không phải [Địa điểm] hay [Dịch vụ] (ví dụ: một ""ghi chú"" cá nhân).\\
Nút ""Đề xuất thông tin"" bị vô hiệu hóa hoặc hệ thống hiển thị thông báo ""Không thể áp dụng cho loại mục này.""
\end{tabular} \\
\bottomrule
\end{longtable}

\renewcommand{\arraystretch}{1.2}
\begin{longtable}{@{}p{4cm} p{11cm}@{}}
\caption{Mô tả Use Case: Điều chỉnh lộ trình}
\label{usecase-1-04} \\
\toprule
\textbf{Use case ID}       & UC-M1-04 \\ \midrule
\textbf{Tên Use case}      & Điều chỉnh lộ trình \\ \midrule
\textbf{Các tác nhân}      & Traveler \\ \midrule
\textbf{Mô tả}             & Traveler thực hiện các thay đổi trên một lộ trình. Sau đó, họ gửi các thay đổi này như một bộ yêu cầu mới để AI tái tạo ra một lộ trình được tối ưu hóa và hợp lý. \\ \midrule
\textbf{Tiền điều kiện}    & Traveler đang xem chi tiết một lộ trình. \\ \midrule
\textbf{Hậu điều kiện}     & Hệ thống trả về một phiên bản lộ trình mới đã được AI tối ưu hóa dựa trên các tinh chỉnh của người dùng, và lưu nó lại. \\ \midrule
\textbf{Luồng sự kiện chính} &
\begin{tabular}[t]{@{}p{10.5cm}@{}}
1. Traveler đang xem một lộ trình trên màn hình.\\
2. Traveler thực hiện nhiều hành động tinh chỉnh khác nhau:\\
– Kéo thả để thay đổi thứ tự hoạt động.\\
– Nhấn vào một hoạt động và nhập prompt mới (ví dụ: ""đổi thành một hoạt động ẩm thực khác"").\\
– Xóa một hoạt động.\\
3. Khi đã hoàn toàn hài lòng với bản nháp mới trên màn hình, Traveler nhấn nút ""Yêu cầu Cập nhật"" (Request).\\
4. Giao diện gửi toàn bộ trạng thái lộ trình đã chỉnh sửa (bản nháp ""thô"") về hệ thống.\\
5. Hệ thống nhận bản nháp này và xem nó như một bộ ràng buộc mới.\\
6. Dựa trên các ràng buộc này, AI sẽ tạo lại các lộ trình mới, đảm bảo nó được tối ưu về mặt logic, thời gian di chuyển, và chi phí trong khi vẫn tôn trọng tối đa các thay đổi của người dùng.\\
7. Hệ thống gửi lộ trình mới và tối ưu này trở lại cho Traveler.\\
8. Giao diện của Traveler được cập nhật để hiển thị phiên bản lộ trình cuối cùng do AI đề xuất.
\end{tabular} \\ \midrule
\textbf{Luồng rẽ nhánh}    & N/A \\ \midrule
\textbf{Luồng ngoại lệ} &
\begin{tabular}[t]{@{}p{10.5cm}@{}}
\textbf{E1: Yêu cầu của người dùng tạo ra xung đột logic:}\\
Xảy ra tại bước 6 của luồng sự kiện chính:\\
6a . AI phân tích các thay đổi của người dùng và phát hiện ra một mâu thuẫn không thể giải quyết.\\
6b . Hệ thống vẫn tiến hành tạo ra một lộ trình thay thế duy nhất, gần nhất và tương đồng nhất với các tinh chỉnh của người dùng.\\
6c . Hệ thống trả về lộ trình thay thế này, kèm theo một cảnh báo giải thích rõ về điểm không phù hợp và sự thay đổi mà AI đã thực hiện (ví dụ: ""Cảnh báo: Yêu cầu của bạn không khả thi tại Vũng Tàu. AI đã đề xuất một lộ trình thay thế gần nhất."").\\
6d . Giao diện của Traveler sẽ hiển thị lộ trình thay thế cuối cùng này.
\end{tabular} \\
\bottomrule
\end{longtable}

\renewcommand{\arraystretch}{1.2}
\begin{longtable}{@{}p{4cm} p{11cm}@{}}
\caption{Mô tả Use Case: Lưu Lộ trình}
\label{usecase-1-05} \\
\toprule
\textbf{Use case ID}       & UC-M1-05 \\ \midrule
\textbf{Tên Use case}      & Lưu Lộ trình \\ \midrule
\textbf{Các tác nhân}      & Traveler \\ \midrule
\textbf{Mô tả}             & Cho phép Traveler lưu lại trạng thái hiện tại của lộ trình (phiên bản nháp hoặc lộ trình do AI tạo) vào tài khoản cá nhân của họ. Điều này đảm bảo họ có thể truy cập lại lộ trình này sau khi thoát ứng dụng. \\ \midrule
\textbf{Tiền điều kiện}    & 
\begin{tabular}[t]{@{}p{10.5cm}@{}}
1. Traveler đã đăng nhập vào hệ thống.\\
2. Traveler đang ở trong giao diện xem hoặc chỉnh sửa một lộ trình.\\
3. Lộ trình này là lộ trình mới (chưa được lưu) hoặc có thay đổi so với lần lưu cuối cùng.
\end{tabular} \\ \midrule
\textbf{Hậu điều kiện}     & (Thành công): Phiên bản hiện tại của lộ trình được lưu trữ vĩnh viễn vào tài khoản của Traveler.\\
(Thất bại): Lộ trình không được lưu, hệ thống hiển thị thông báo lỗi. \\ \midrule
\textbf{Luồng sự kiện chính} &
\begin{tabular}[t]{@{}p{10.5cm}@{}}
1. Traveler đang xem hoặc chỉnh sửa một lộ trình mà họ ưng ý.\\
2. Traveler nhấn nút ""Lưu Lộ trình"".\\
3. Hệ thống kiểm tra xem đây là lần lưu đầu tiên cho lộ trình này hay là cập nhật (ghi đè) một lộ trình đã có.\\
4. (Trường hợp lưu lần đầu): Hệ thống hiển thị cửa sổ yêu cầu Traveler nhập ""Tên Lộ trình"".\\
5. (Trường hợp lưu lần đầu): Traveler nhập tên và nhấn ""Xác nhận"".\\
6. Hệ thống thu thập toàn bộ dữ liệu của lộ trình hiện tại (danh sách hoạt động, ngày giờ, địa điểm, ghi chú...).\\
7. Hệ thống lưu/cập nhật dữ liệu lộ trình này vào cơ sở dữ liệu, liên kết với ID của Traveler.\\
8. Hệ thống hiển thị thông báo ""Đã lưu lộ trình thành công!"".
\end{tabular} \\ \midrule
\textbf{Luồng rẽ nhánh}    & 
\begin{tabular}[t]{@{}p{10.5cm}@{}}
\textbf{A1: Cập nhật (ghi đè) lộ trình đã có:}\\
Tại bước 3, hệ thống phát hiện lộ trình này đã được lưu trước đó.\\
Hệ thống bỏ qua bước 4 và 5.\\
Luồng sự kiện tiếp tục từ bước 6 (thu thập dữ liệu và ghi đè lên phiên bản cũ).
\end{tabular} \\ \midrule
\textbf{Luồng ngoại lệ} &
\begin{tabular}[t]{@{}p{10.5cm}@{}}
\textbf{E1: Lỗi kết nối hoặc Cơ sở dữ liệu:}\\
Tại bước 7, hệ thống không thể lưu dữ liệu do lỗi máy chủ hoặc mất kết nối.\\
Hệ thống hiển thị thông báo ""Lưu thất bại, vui lòng thử lại.""\\[3pt]
\textbf{E2: Tên lộ trình không hợp lệ (khi lưu lần đầu):}\\
Tại bước 5, Traveler nhập tên đã tồn tại hoặc chứa ký tự không hợp lệ.\\
Hệ thống chặn việc lưu và hiển thị thông báo ""Tên lộ trình đã tồn tại (hoặc không hợp lệ). Vui lòng chọn tên khác.""
\end{tabular} \\
\bottomrule
\end{longtable}

\renewcommand{\arraystretch}{1.2}
\begin{longtable}{@{}p{4cm} p{11cm}@{}}
\caption{Mô tả Use Case: Xóa Lộ trình}
\label{usecase-1-06} \\
\toprule
\textbf{Use case ID}       & UC-M1-06 \\ \midrule
\textbf{Tên Use case}      & Xóa Lộ trình \\ \midrule
\textbf{Các tác nhân}      & Traveler \\ \midrule
\textbf{Mô tả}             & Cho phép Traveler xóa vĩnh viễn một lộ trình đã lưu ra khỏi tài khoản của họ. \\ \midrule
\textbf{Tiền điều kiện}    & 
\begin{tabular}[t]{@{}p{10.5cm}@{}}
1. Traveler đã đăng nhập vào hệ thống.\\
2. Traveler đang ở trang quản lý danh sách các lộ trình đã lưu của mình.
\end{tabular} \\ \midrule
\textbf{Hậu điều kiện}     & Lộ trình được chọn sẽ bị xóa vĩnh viễn khỏi tài khoản của Traveler và không còn hiển thị trong danh sách. \\ \midrule
\textbf{Luồng sự kiện chính} &
\begin{tabular}[t]{@{}p{10.5cm}@{}}
1. Traveler tìm đến lộ trình mà họ muốn xóa trong danh sách các lộ trình đã lưu.\\
2. Traveler nhấn vào nút hoặc biểu tượng ""Xóa"" tương ứng với lộ trình đó.\\
3. Hệ thống hiển thị một hộp thoại xác nhận để tránh xóa nhầm, ví dụ: ""Bạn có chắc chắn muốn xóa vĩnh viễn lộ trình '[Tên lộ trình]' không? Hành động này không thể hoàn tác.""\\
4. Traveler nhấn nút ""Xác nhận Xóa"".\\
5. Hệ thống gửi yêu cầu xóa đến máy chủ.\\
6. Máy chủ xóa vĩnh viễn bản ghi của lộ trình đó khỏi cơ sở dữ liệu.\\
7. Hệ thống làm mới lại danh sách lộ trình trên giao diện, lộ trình vừa bị xóa sẽ biến mất.\\
8. Hệ thống hiển thị một thông báo ngắn gọn ""Đã xóa lộ trình thành công.""
\end{tabular} \\ \midrule
\textbf{Luồng rẽ nhánh}    & 
\begin{tabular}[t]{@{}p{10.5cm}@{}}
\textbf{A1: Hủy bỏ hành động xóa:}\\
Thời điểm: Xảy ra tại bước 4 của Luồng sự kiện chính.\\
Mô tả: Traveler nhấn nút ""Hủy"" trong hộp thoại xác nhận. Hệ thống sẽ đóng hộp thoại và không có hành động nào được thực hiện. Use case kết thúc.
\end{tabular} \\ \midrule
\textbf{Luồng ngoại lệ} &
\begin{tabular}[t]{@{}p{10.5cm}@{}}
\textbf{E1: Lỗi kết nối hoặc lỗi máy chủ:}\\
Thời điểm: Xảy ra tại bước 6.\\
Mô tả: Hệ thống không thể xóa lộ trình do lỗi kết nối hoặc lỗi phía máy chủ. Hệ thống sẽ hiển thị thông báo lỗi (ví dụ: ""Xóa thất bại, vui lòng thử lại."") và lộ trình vẫn còn trong danh sách.\\[4pt]
\textbf{E2: Lộ trình không còn tồn tại:}\\
Thời điểm: Xảy ra tại bước 5.\\
Mô tả: Lộ trình đã bị xóa từ một thiết bị hoặc phiên làm việc khác trước khi yêu cầu này được thực hiện. Hệ thống sẽ hiển thị thông báo ""Lộ trình này không còn tồn tại."" và tự động làm mới danh sách.
\end{tabular} \\
\bottomrule
\end{longtable}
